\section{Experimental Setup}
\label{sec:experimental-setup}

\subsection{Evaluation Objectives}

The experiments were designed to evaluate the complete pipeline from
retrieval to the final answer while separately assessing the quality of each
component. Five research questions were formulated as follows:

\begin{description}
    \item[RQ1 --- Retrieval quality.]
    Does the system retrieve the rules, events, and articles relevant to the
    question?

    \item[RQ2 --- Causal-path quality.]
    Does the system recover the correct ordering of events and rules along the
    gold causal path? If the gold path appears in the candidate set, does the
    system rank it first?

    \item[RQ3 --- Verification quality.]
    Does the query-aware verifier correctly classify supported claims and
    direct-edge claims that should be rejected?

    \item[RQ4 --- Answer quality.]
    Does the answer generator return the requested type of information,
    preserve the content of the gold answer, and cite the correct legal
    articles?

    \item[RQ5 --- Effect of the verification engine.]
    How do Structural SCM and node-deletion path ablation differ in output
    quality, processing time, and trace detail?
\end{description}

In addition to aggregate results, the study analyzes metrics by question type,
difficulty, and counterfactual requirement to identify the principal error
groups.

\subsection{Evaluation Protocol}

The experiments use the complete benchmark of \(250\) questions. No training
or fine-tuning is performed on the benchmark. The system receives only the
question as input at inference time; the evaluation fields are reserved for
Step 7.

No train--validation--test split is used because the principal pipeline
components are not trained on the benchmark. However, because the benchmark is
generated from the same graph used by the retriever, the results constitute a
closed-corpus evaluation rather than an assessment of out-of-domain
generalization.

Each configuration executes the following steps in sequence:

\begin{enumerate}
    \item retrieve events, rules, and causal paths;
    \item verify the paths and claim;
    \item generate an extractive answer;
    \item normalize the prediction to the common schema;
    \item compute metrics using Step 7.
\end{enumerate}

The evaluator is run with the \texttt{strict} option. Evaluation therefore
stops if any prediction is missing or if a prediction contains an ID not
present in the benchmark. Both paired runs must contain all \(250\) predictions
and have no pipeline errors before they are included in the comparison.

\subsection{Component Versions}

The versions used in the principal experiments are presented in
Table~\ref{tab:versions}.

\begin{table}[t]
    \centering
    \caption{Versions of the components in the experimental pipeline.}
    \label{tab:versions}
    \begin{tabular}{ll}
        \hline
        \textbf{Component} & \textbf{Version} \\
        \hline
        Benchmark & \texttt{2.0} \\
        Step 4 verifier & \texttt{4.1-query-aware-legal-scm} \\
        Step 5 answer generator & \texttt{5.1-query-aware-grounded-answer} \\
        Batch runner & \texttt{2.1-query-aware-answer} \\
        Step 7 evaluator & \texttt{2.0-current-pipeline-schema} \\
        \hline
    \end{tabular}
\end{table}

Version information is stored directly in the artifact metadata. A prediction
without the correct version provenance is not considered part of the principal
paired experiment.

\subsection{Execution Environment}

The two paired runs were executed in the same server environment and used the
same resources. To ensure reproducibility, the final version of the paper must
include complete hardware and software information in
Table~\ref{tab:environment}.

\begin{table}[t]
    \centering
    \caption{Execution environment. The hardware fields must be completed
    using information from the server on which the experiments were run before
    publication.}
    \label{tab:environment}
    \begin{tabular}{ll}
        \hline
        \textbf{Component} & \textbf{Information} \\
        \hline
        Operating system & \textit{[To be completed from server]} \\
        Python version & \textit{[To be completed from server]} \\
        CPU & \textit{[To be completed from server]} \\
        RAM & \textit{[To be completed from server]} \\
        GPU, if used & \textit{[To be completed from server]} \\
        Embedding model & \texttt{BAAI/bge-m3} \\
        Vector index & FAISS \\
        Answer provider & Extractive \\
        \hline
    \end{tabular}
\end{table}

BGE-M3 is used for semantic retrieval and semantic mapping
\cite{chen2024bgem3}. FAISS is used for nearest-neighbor search over the
causal memory \cite{johnson2019billion}.

CPU, RAM, GPU, and library-version information is not inferred from the
artifacts. These fields must be obtained directly from the experimental
server. This is particularly important when reporting runtime.

\subsection{Retrieval Configuration}

The two paired runs use the same embeddings, graph, causal memory, and FAISS
index. The Step 3 hyperparameters are presented in
Table~\ref{tab:retrieval-hyperparameters}.

\begin{table}[t]
    \centering
    \caption{Hyperparameters of the multi-hop causal retriever.}
    \label{tab:retrieval-hyperparameters}
    \begin{tabular}{lr}
        \hline
        \textbf{Parameter} & \textbf{Value} \\
        \hline
        \texttt{event\_top\_k} & \(8\) \\
        \texttt{direct\_rule\_top\_k} & \(8\) \\
        \texttt{semantic\_pool\_size} & \(100\) \\
        \texttt{max\_hops} & \(2\) \\
        \texttt{max\_paths\_per\_event} & \(30\) \\
        \texttt{max\_candidate\_rules} & \(200\) \\
        \texttt{final\_top\_k} & \(12\) \\
        \texttt{min\_event\_score} & \(0.20\) \\
        \texttt{min\_rule\_score} & \(0.15\) \\
        \texttt{direction} & \texttt{both} \\
        \hline
    \end{tabular}
\end{table}

The \texttt{direction=both} option allows the retriever to traverse the graph
in both the forward and reverse directions. This is necessary because the
benchmark includes both forward multihop and reverse-reasoning questions.
However, bidirectional traversal does not guarantee that paths will be ranked
appropriately for the answer intent; this effect is analyzed in
Section~\ref{sec:discussion}.

\subsection{Query-Aware Verification Configuration}

The shared Step 4 hyperparameters are presented in
Table~\ref{tab:verification-hyperparameters}.

\begin{table}[t]
    \centering
    \caption{Hyperparameters of query-aware verification.}
    \label{tab:verification-hyperparameters}
    \begin{tabular}{lr}
        \hline
        \textbf{Parameter} & \textbf{Value} \\
        \hline
        \texttt{cf\_top\_k} & \(5\) \\
        \texttt{mapping\_top\_k} & \(5\) \\
        \texttt{mapping\_threshold} & \(0.42\) \\
        \texttt{max\_cf\_hops} & \(3\) \\
        \texttt{max\_cf\_paths} & \(30\) \\
        \texttt{verified\_top\_k} & \(10\) \\
        \texttt{keep\_threshold} & \(0.52\) \\
        \texttt{reject\_threshold} & \(0.34\) \\
        Alternative path threshold & \(0.35\) \\
        Semantic mapping & Enabled \\
        \hline
    \end{tabular}
\end{table}

Semantic mapping is enabled in both runs. Thus, the difference between
Structural SCM and path ablation does not arise because one configuration uses
embedding-based mapping while the other does not.

\subsection{Two Verification Modes}

The paired experiment comprises two configurations.

\subsubsection{Structural SCM}

The first configuration uses:

\begin{verbatim}
counterfactual_mode = structural_scm
\end{verbatim}

LegalSCM performs:

\begin{itemize}
    \item factual path validation;
    \item rule-coverage checking;
    \item three-valued fixed-point inference over the normative encoding;
    \item hard intervention on the mediator;
    \item renewed inference of the counterfactual outcome;
    \item recording of disabled rules, newly activated rules, and the
    intervention trace.
\end{itemize}

The node-deletion baseline is retained in the trace as a diagnostic, but the
final structural result is prioritized when LegalSCM completes successfully.

\subsubsection{Path Ablation}

The second configuration uses:

\begin{verbatim}
counterfactual_mode = path_ablation
\end{verbatim}

The mediator is removed from the graph, and the system searches for a bounded
alternative path from the seed event to the outcome. This configuration does
not perform structural fixed-point inference and is not interpreted as a
\(do(\cdot)\) intervention.

\subsubsection{Conditions for Paired Comparison}

The two runs use the same:

\begin{itemize}
    \item benchmark of \(250\) questions;
    \item retriever and retrieval-output schema;
    \item embedding model and semantic mapping;
    \item graph, rule corpus, memory, and index;
    \item verification thresholds;
    \item Step 5 answer generator;
    \item extractive provider;
    \item evaluator and metrics.
\end{itemize}

The only independent variable is \texttt{counterfactual\_mode}. Therefore,
the difference between the two runs can be attributed to the verification
engine within the scope of the current implementation.

\subsection{Answer-Generation Configuration}

The Step 5 hyperparameters are presented in
Table~\ref{tab:generation-hyperparameters}.

\begin{table}[t]
    \centering
    \caption{Hyperparameters of answer generation.}
    \label{tab:generation-hyperparameters}
    \begin{tabular}{lr}
        \hline
        \textbf{Parameter} & \textbf{Value} \\
        \hline
        Provider & \texttt{extractive} \\
        \texttt{max\_evidence} & \(8\) \\
        \texttt{answer\_max\_paths} & \(6\) \\
        \texttt{max\_context\_chars} & \(18\,000\) \\
        \texttt{max\_tokens} & \(1\,200\) \\
        \texttt{temperature} & \(0.1\) \\
        \texttt{timeout} & \(180\) seconds \\
        \texttt{min\_verification\_score} & \(0.45\) \\
        Include uncertain evidence & No \\
        Extractive fallback & Enabled \\
        \hline
    \end{tabular}
\end{table}

Because the primary provider is extractive, the \texttt{max\_tokens} and
\texttt{temperature} parameters do not directly affect answer content in the
paired experiment. They are retained in the configuration so that the same
schema can support other generative providers.

Neither an LLM judge nor LLM generation is used in the principal results. This
choice reduces stochastic variation and enables direct analysis of the effects
of retrieval, verification, and answer intent.

\subsection{Evaluator Configuration}

Step 7 is run with the following values:

\[
k \in \{1,3,5,10\}.
\]

The evaluator uses:

\begin{itemize}
    \item all \(250\) gold samples;
    \item all \(250\) predictions;
    \item \texttt{strict} mode;
    \item the same metric implementation for both configurations;
    \item separate output directories for Structural SCM and path ablation.
\end{itemize}

The principal result directories are:

\begin{verbatim}
evaluation_results/structural_scm_v41_step5_v51/
evaluation_results/path_ablation_v41_step5_v51/
\end{verbatim}

Older reports without the \texttt{step5\_v51} suffix are not used in the
paired comparison.

\subsection{Retrieval Metrics}

Let \(\mathcal{G}\) denote the set of gold IDs and
\(\widehat{\mathcal{R}}_k\) the top \(k\) ranked IDs. Precision and Recall at
\(k\) are computed as:

\begin{equation}
\mathrm{Precision@}k
=
\frac{
    \left|
        \widehat{\mathcal{R}}_k
        \cap
        \mathcal{G}
    \right|
}{
    k
},
\end{equation}

\begin{equation}
\mathrm{Recall@}k
=
\frac{
    \left|
        \widehat{\mathcal{R}}_k
        \cap
        \mathcal{G}
    \right|
}{
    \left|\mathcal{G}\right|
}.
\end{equation}

Hit@\(k\) is assigned a value of \(1\) if at least one gold item appears in
the top \(k\), and \(0\) otherwise.

The metrics are computed separately for:

\begin{itemize}
    \item rule IDs;
    \item event IDs;
    \item article IDs.
\end{itemize}

The evaluator also reports:

\begin{itemize}
    \item Macro Precision, Recall, and F1;
    \item Micro Precision, Recall, and F1;
    \item Exact Set Match;
    \item Mean Reciprocal Rank (MRR);
    \item Mean Average Precision (MAP).
\end{itemize}

In the paper, Rule Recall@5 and Event Recall@5 are used as the principal
retrieval metrics because a benchmark sample may contain multiple gold rules
or events along the same path.

\subsection{Causal-Path Metrics}

The gold path and predicted path are normalized into sequences of directed
edges:

\[
\left[
(v_0,v_1),
(v_1,v_2),
\ldots,
(v_{h-1},v_h)
\right].
\]

\paragraph{Top-1 Exact Path Match.}

This metric is assigned a value of \(1\) if the path selected by the pipeline
as the primary path exactly matches the gold path in edge order, and \(0\)
otherwise.

\paragraph{Oracle Exact Path Match.}

The evaluator examines the candidate paths returned by the retriever and
selects the path that best matches the gold path. The oracle metric is used
only to diagnose retrieval capability; it does not represent the path actually
used by the system to generate the answer.

The gap:

\[
\mathrm{OracleExact}
-
\mathrm{Top1Exact}
\]

is used to identify path-ranking errors. If both the oracle and top-1 results
are incorrect, the error is more likely to lie in path retrieval.

\paragraph{Edge-level metrics.}

Predicted edges and gold edges are treated as two sets for computing precision,
recall, and F1.

\paragraph{Event-level metrics.}

The set of events appearing on the predicted path is compared with the set of
events on the gold path.

\paragraph{Final Event Accuracy.}

This metric checks whether the final outcome event of the predicted path
matches the gold final event.

\paragraph{Path Rule F1.}

The set of rules supporting the predicted path is compared with the gold rules
for the path.

For branch and convergence questions, a linear exact-path metric does not fully
represent the structure that must be predicted. Their path metrics are
therefore reported but must be interpreted jointly with answer and citation
metrics.

\subsection{Verification Metrics}

The verification decision has three labels:

\[
\mathcal{D}
=
\{
\texttt{SUPPORTED},
\texttt{REJECT\_DIRECT\_CLAIM},
\texttt{UNCERTAIN}
\}.
\]

Accuracy is computed as:

\begin{equation}
\mathrm{Accuracy}
=
\frac{
    \text{number of predictions with the correct decision}
}{
    \text{total number of predictions}
}.
\end{equation}

Because the benchmark has imbalanced labels, the evaluator also reports:

\begin{itemize}
    \item per-class precision, recall, and F1;
    \item Macro-F1 over classes with gold support;
    \item Balanced Accuracy;
    \item a confusion matrix.
\end{itemize}

Balanced Accuracy is the mean recall over classes represented in the gold
data:

\begin{equation}
\mathrm{BalancedAccuracy}
=
\frac{1}{|\mathcal{D}^{+}|}
\sum_{d \in \mathcal{D}^{+}}
\mathrm{Recall}_d,
\end{equation}

where \(\mathcal{D}^{+}\) is the set of classes with at least one gold
sample.

Counterfactual verification accuracy is reported separately for samples with
\texttt{requires\_counterfactual=true}. However, as discussed in
Section~\ref{sec:data}, this subset primarily evaluates negative direct-edge
claims.

\subsection{Answer Metrics}

Before evaluation, the predicted answer and gold answer are:

\begin{itemize}
    \item converted to lowercase;
    \item normalized for whitespace;
    \item stripped of unnecessary punctuation;
    \item tokenized;
    \item retained in Vietnamese.
\end{itemize}

Token Precision and Token Recall are computed from the number of overlapping
tokens between the prediction and the gold answer:

\begin{equation}
P_{\mathrm{token}}
=
\frac{
    |\widehat{T} \cap T|
}{
    |\widehat{T}|
},
\end{equation}

\begin{equation}
R_{\mathrm{token}}
=
\frac{
    |\widehat{T} \cap T|
}{
    |T|
},
\end{equation}

where \(\widehat{T}\) and \(T\) are the token multisets of the predicted and
gold answers, respectively.

Answer Token F1 is computed as:

\begin{equation}
F1_{\mathrm{token}}
=
\frac{
    2P_{\mathrm{token}}R_{\mathrm{token}}
}{
    P_{\mathrm{token}}+R_{\mathrm{token}}
}.
\end{equation}

If both precision and recall are \(0\), F1 is set to \(0\).

ROUGE-L F1 is computed from the longest common subsequence between the
prediction and the gold answer \cite{lin2004rouge}. This metric retains some
information about token order that bag-of-token F1 does not capture.

The evaluator also reports Normalized Exact Match. However, because extractive
answers often include an additional causal chain and citations, Exact Match is
not selected as a principal metric.

\subsection{Citation Metrics}

Citations are normalized to article IDs. For example:

\[
\text{``Article 15''}
\longrightarrow
\texttt{15}.
\]

Duplicate citations are removed before evaluation. Given the set of predicted
citations \(\widehat{\mathcal{C}}\) and the set of gold citations
\(\mathcal{C}\), the evaluator computes:

\begin{equation}
P_{\mathrm{citation}}
=
\frac{
    |
    \widehat{\mathcal{C}}
    \cap
    \mathcal{C}
    |
}{
    |
    \widehat{\mathcal{C}}
    |
},
\end{equation}

\begin{equation}
R_{\mathrm{citation}}
=
\frac{
    |
    \widehat{\mathcal{C}}
    \cap
    \mathcal{C}
    |
}{
    |
    \mathcal{C}
    |
},
\end{equation}

and the corresponding harmonic F1. Citation Exact Match is assigned a value
of \(1\) when the two sets of article IDs are identical.

Citation F1 evaluates whether the correct legal articles are selected, but it
does not by itself establish that the answer's explanation is consistent with
the complete content of those articles.

\subsection{Runtime Measurement}

The batch runner records the end-to-end runtime for each sample, including
retrieval, verification, and answer generation. Step 7 reports:

\begin{itemize}
    \item average runtime per sample;
    \item median runtime per sample.
\end{itemize}

The run log also records the wall-clock time of the complete batch. The two
types of runtime are reported separately:

\begin{itemize}
    \item per-sample runtime reflects the processing time recorded for each
    question;
    \item batch wall-clock time additionally includes model initialization,
    resource loading, checkpointing, and artifact writing.
\end{itemize}

Because the hardware information has not yet been fully populated, runtime is
used only for relative comparison between two configurations executed on the
same server and is not treated as a speed benchmark that generalizes to other
environments.

\subsection{Error Analysis}

Step 7 assigns one or more error types to each sample. The principal error
groups are:

\begin{description}
    \item[\texttt{RULE\_RETRIEVAL\_MISS}]
    The gold rules are not fully present in the top-\(k\) results.

    \item[\texttt{EVENT\_RETRIEVAL\_MISS}]
    The gold events are not fully retrieved.

    \item[\texttt{TOP1\_PATH\_MISMATCH}]
    The primary path does not match the gold path.

    \item[\texttt{PATH\_RANKING\_ERROR}]
    The gold path appears in the candidate set but is not ranked first.

    \item[\texttt{PATH\_RETRIEVAL\_ERROR}]
    The gold path does not appear in the candidate set.

    \item[\texttt{VERIFICATION\_ERROR}]
    The predicted decision differs from the gold decision.

    \item[\texttt{LOW\_ANSWER\_TOKEN\_F1}]
    The answer's token overlap is below the error-analysis threshold.

    \item[\texttt{CITATION\_MISS}]
    At least one gold citation is missing.

    \item[\texttt{EXTRA\_CITATION}]
    The answer contains a citation not included in the gold set.

    \item[\texttt{PIPELINE\_ERROR}]
    The sample does not produce a complete prediction.
\end{description}

Error types are used only for post-inference analysis and are not used to
adjust predictions.

\subsection{Subgroup Analysis}

Metrics are aggregated in three ways:

\begin{enumerate}
    \item by question type;
    \item by difficulty;
    \item by the \texttt{requires\_counterfactual} flag.
\end{enumerate}

Analysis by question type identifies the system's strengths and weaknesses on
forward, reverse, bridge, chain, branch, and convergence questions. Analysis
by difficulty measures degradation on samples with many candidate rules or
complex structures. Analysis of the counterfactual subset separately evaluates
the direct claims that should be rejected.

\subsection{Principles for Interpreting the Paired Ablation}

The paired experiment is interpreted according to the following principles:

\begin{enumerate}
    \item If retrieval and path metrics are identical, this is expected because
    the verification engine is applied after retrieval.

    \item If verification or answer metrics differ, the difference may be
    related to the final decision, primary path, or evidence retained by Step
    4.

    \item If quality metrics are identical but runtime differs, no conclusion
    should be drawn that Structural SCM improves accuracy.

    \item A more detailed trace is an auditability benefit, not a quality
    improvement unless it is separately evaluated through a metric or expert
    study.

    \item Oracle Exact Path must not be substituted for Top-1 Exact Path when
    reporting deployable performance.

    \item Counterfactual accuracy on \(20\) negative direct-edge samples must
    not be treated as comprehensive validation of hard-intervention reasoning.
\end{enumerate}

Because the two configurations are deterministic and were each run only once
on the same benchmark, the study reports a descriptive paired comparison. No
statistical significance test is performed, and no standard deviation across
multiple random seeds is reported. When the two configurations produce
identical results, the outcome is described as a \emph{null quality
ablation}, not as evidence of general equivalence across all datasets.
